% Chapter: Reinforcement Learning Algorithms
% Two sections: Classical Synthesis (Sutton-Baird), Deep Learning Era (Mnih-AlphaGo)

%%%%%%%%%%%%%%%%%%%%%%%%%%%%%%%%%%%%%%%%%%%%%%%%%%%%%%%%%%%%%%%%%%%%%%%%%%%%%%%
\subsection{The Classical Synthesis}
\label{sec:classical_synthesis}
%%%%%%%%%%%%%%%%%%%%%%%%%%%%%%%%%%%%%%%%%%%%%%%%%%%%%%%%%%%%%%%%%%%%%%%%%%%%%%%

\subsubsection{Monte Carlo Estimation}

When $P(s'|s,a)$ and $r(s,a)$ are unknown, the obvious approach is to estimate value functions directly from sampled \textit{episodes} $(s_0, a_0, r_1, s_1, a_1, r_2, \ldots, s_T)$. The realized \textit{return} from state $s_t$ is
\begin{equation}
G_t = \sum_{k=0}^{T-t-1} \gamma^k r_{t+k+1}
\end{equation}
First-visit MC prediction averages $G_t$ over episodes for each state $s$, counting only its first occurrence per episode. Each first-visit return is an independent draw from the return distribution, so the sample mean converges almost surely to $V^\pi(s)$ by the strong law of large numbers \citep{sutton2018}. An incremental update is, $$V(s) \leftarrow V(s) + \alpha[G_t - V(s)]$$

For MC control, one estimates action-values $Q(s,a)$ by averaging first-visit returns from each state-action pair, then improve the policy greedily: $\pi(s) = \argmax_a Q(s,a)$. Under exploring starts (every $(s,a)$ pair begins an episode infinitely often), this alternation converges to $Q^*$. \footnote{\citet{tsitsiklis2002} proved convergence even when the policy improves after every episode rather than waiting for complete evaluation. In practice, exploring starts is infeasible, so on-policy variants use $\varepsilon$-greedy exploration instead. These converge to $Q^*$ provided the exploration schedule satisfies the greedy-in-the-limit with infinite exploration (GLIE) condition: every state-action pair is visited infinitely often, and $\varepsilon_t \to 0$ so the policy converges to greedy in the limit.}

\subsubsection{Sutton (1988)}

Monte Carlo has two limitations. The agent must wait for episode termination to compute $G_t$, ruling out continuing tasks. And $G_t$ is unbiased ($\mathbb{E}[G_t \mid S_t = s] = V^\pi(s)$) but high-variance, because it sums random rewards over the entire \textit{trajectory}.

\citet{sutton1988} proposed temporal difference (TD) learning to fix both problems. TD(0) replaces the full return $G_t$ with a one-step target:
\begin{equation}
V(s_t) \leftarrow V(s_t) + \alpha \left[ r_{t+1} + \gamma V(s_{t+1}) - V(s_t) \right]
\end{equation}
The target $r_{t+1} + \gamma V(s_{t+1})$ depends on one random reward and one random transition, so its variance is low. The cost is bias. The bootstrap target $V(s_{t+1})$ is the agent's current estimate, not the true value. This is ``\textit{bootstrapping}''.\footnote{See Section~\ref{section:language} for the distinction between RL bootstrapping and Efron's resampling procedure.} As $V$ improves, the bias shrinks; the low variance persists regardless. Sutton demonstrated this tradeoff on a five-state random walk where TD(0) converged faster than Monte Carlo with less data.

The general TD($\lambda$) update interpolates between these extremes through an eligibility trace.\footnote{The eligibility trace records which states were recently visited, allowing credit assignment to propagate backward in time. States visited more recently receive stronger updates when the TD error is observed.}
\begin{equation}
V(s_t) \leftarrow V(s_t) + \alpha \, \delta_t \, e_t(s)
\end{equation}
where $\delta_t = r_{t+1} + \gamma V(s_{t+1}) - V(s_t)$ is the TD error and $e_t(s) = \gamma \lambda \, e_{t-1}(s) + \mathbbm{1}\{s = s_t\}$ is the eligibility trace for state $s$. Setting $\lambda = 0$ yields TD(0); setting $\lambda = 1$ recovers Monte Carlo returns. Intermediate $\lambda$ trades off variance against bias. \footnote{\citet{dayan1992} proved convergence of TD($\lambda$) for general $\lambda$ in the tabular case; \citet{jaakkola1994} gave a unified stochastic approximation proof covering both TD and Q-learning; \citet{tsitsiklis1997} extended the analysis to linear function approximation.}

\subsubsection{Watkins (1989)}

TD(0) learns value functions $V(s)$, but converting these to actions (the control problem) still requires knowing transition probabilities. Given $V(s')$ for all successor states, the agent needs to know which action leads to which successor. 

\citet{WatkinsDayan1992}, formalizing Watkins's 1989 PhD thesis, provided the solution. Instead of learning $V(s')$ learn $Q(s,a)$ (the action value function or the ``quality" of actions function) directly, the expected return from taking action $a$ in state $s$ and then behaving optimally. The optimal policy is then $\pi^*(s) = \argmax_a Q^*(s,a)$, requiring no model to act. The Bellman optimality equation provides the fixed-point condition:
\begin{equation}
Q^*(s,a) = \mathbb{E}\left[r + \gamma \max_{a'} Q^*(s',a') \mid s, a\right]
\end{equation}
Q-learning achieves this via the update
\begin{equation}
Q(s_t, a_t) \leftarrow Q(s_t, a_t) + \alpha \left[ r_{t+1} + \gamma \max_{a'} Q(s_{t+1}, a') - Q(s_t, a_t) \right]
\end{equation}
Q-learning converges to $Q^*$ under standard regularity conditions (Section~\ref{sec:stochastic_approx}). The maximization over $a'$ makes Q-learning \textit{off-policy}:\footnote{See Section~\ref{section:language} for the on-policy/off-policy distinction. Q-learning learns about the greedy policy $\pi^*(s) = \argmax_a Q(s,a)$ while collecting data with an exploratory policy. This exploratory policy needs only to be sufficiently ``exploratory'' and admits a wide range of policies; including a fully random policy.} the update target uses the greedy action at the next state regardless of the action actually taken. Therefore the agent can follow an $\varepsilon$-greedy exploration strategy or even a fully random policy, while learning about the optimal policy directly.

\subsubsection{Williams (1992)}

Value-based methods learn an action-value function and derive a policy from it. \citet{williams1992} derived an alternative, policy gradients, that optimize the policy directly by gradient ascent on expected return
\begin{equation}
\nabla_\theta J(\theta) = \mathbb{E}_{\pi_\theta} \left[ \sum_{t=0}^{T} \nabla_\theta \log \pi_\theta(a_t|s_t) \, G_t \right]
\end{equation}
where $G_t = \sum_{k=0}^{T-t} \gamma^k r_{t+k+1}$ is the discounted return from time $t$. The log-derivative trick allows the gradient to be estimated from sampled trajectories without differentiating through the environment dynamics.

When the action space is continuous, Q-learning requires computing $\max_a Q(s,a)$ at every update, a nested optimization problem.\footnote{In continuous action spaces, $\max_a Q(s,a)$ has no closed-form solution in general and must be solved numerically at every Bellman update. Discretizing a $d$-dimensional action space on a grid of $m$ points per dimension costs $O(m^d)$ evaluations per update step.} Policy gradients sidestep the issue. Parameterize the policy as a Gaussian $\pi_\theta(a|s) = \mathcal{N}(\mu_\theta(s),\, \sigma_\theta^2(s))$, sample an action, observe the return, and update $\theta$ by REINFORCE. Continuous-control results in reinforcement learning descend from this framework.

%REINFORCE has a few practical issues, which motivated. First, because $G_t$ is a sum of stochastic rewards, gradient estimates can be noisy. Variance reduction techniques, including baselines $b(s_t)$ subtracted from $G_t$, are essential. One approach is to also estimate $V(s_t)$ and use it as a baseline. Secondly, REINFORCE uses each trajectory exactly once and discards it making it very sample-inefficient. This led to parallelization to gather more data faster and importance-sampling methods that re-use older data. Thirdly, a few outliers could destabilize learning if the gradient update is too severe. This led to KL trust regions and other limits on policy updates.  Thirdly, there is no mechanism to prevent the policy from collapsing into a deterministic policy early in training. This led to entropy regularization to encourage exploration by penalizing low-entropy policies.

\subsubsection{Tesauro (1994)}

\citet{tesauro1994}'s TD-Gammon demonstrated that temporal difference learning with neural network function approximation could achieve expert-level play in backgammon, a domain with approximately $10^{20}$ legal positions. A feedforward neural network estimated the probability of winning from any board position, trained by self-play using TD($\lambda$) with $\lambda = 0.7$. At each turn, the current player selected the legal move maximizing the network's value estimate. As the network improved, it generated stronger opponents for itself, producing harder training games that drove further improvement. Backgammon's dice rolls provided natural exploration without requiring an explicit mechanism.\footnote{Version 2.1, trained on 1,500,000 self-play games, achieved near-parity with former world champion Bill Robertie. The logistic sigmoid activation $\sigma(x) = 1/(1+e^{-x})$ used in the network is identical to the binary logit link function in discrete choice.}

\subsubsection{SARSA (1994)}

Q-learning learns the optimal action-value function regardless of the policy generating experience. This off-policy property is useful but introduces complications when combined with function approximation. \citet{rummery1994} introduced SARSA as an \textit{on-policy}\footnote{See Section~\ref{section:language} for the on-policy/off-policy distinction.} alternative that learns the value of the policy actually being followed. The name derives from the quintuple $(s_t, a_t, r_{t+1}, s_{t+1}, a_{t+1})$ used in each update:
\begin{equation}
Q(s_t, a_t) \leftarrow Q(s_t, a_t) + \alpha \left[ r_{t+1} + \gamma Q(s_{t+1}, a_{t+1}) - Q(s_t, a_t) \right]
\end{equation}
The key difference from Q-learning is that SARSA bootstraps from the action $a_{t+1}$ actually taken at the next state, rather than the greedy action $\argmax_{a'} Q(s_{t+1}, a')$. This makes the algorithm on-policy, since the target depends on the behavior policy generating the data.

If the agent follows an $\varepsilon$-greedy policy, SARSA converges to $Q^{\varepsilon\text{-greedy}}$, not $Q^*$\footnote{SARSA converges to $Q^*$ under GLIE (greedy in the limit with infinite exploration). A GLIE schedule explores all state-action pairs infinitely often but converges to the greedy policy asymptotically. The standard $\varepsilon$-greedy policy with $\varepsilon_t \to 0$ is one such schedule.} policies and standard step-size conditions (Section~\ref{sec:stochastic_approx}). This distinction matters when exploration is costly. Consider the cliff-walking problem, where an agent must traverse a gridworld with a cliff along one edge. The optimal path runs along the cliff edge (shortest route), but the $\varepsilon$-greedy policy occasionally falls off. Q-learning learns the optimal path because it evaluates the greedy policy; the agent falls off during learning but the Q-values reflect the optimal route. SARSA learns a safer path further from the cliff because it evaluates the actual exploratory policy; it accounts for the fact that exploration sometimes leads to catastrophic states. 

\subsubsection{Baird (1995)}

\citet{Baird1995} constructed a six-state star MDP demonstrating divergence of Q-learning with linear function approximation. The MDP has five outer states that all transition to a single inner state under the target policy. The off-policy behavior samples states uniformly. With linear function approximation, the weights grow without bound under repeated Q-learning updates. The source of instability is the interaction of three components, namely bootstrapping (updating from estimated values rather than observed returns), off-policy learning (training on data from a different policy than the target), and function approximation (representing the value function with a parameterized model). \citet{sutton2018} later named this the deadly triad. Any two components can be combined safely; all three together permit divergence. The mechanism underlying this instability is analyzed in Section~\ref{sec:deadly_triad}.

This result explains the asymmetry between Tesauro's success and Baird's failure. TD-Gammon used bootstrapping and function approximation but was on-policy, with training data coming from the same self-play policy whose value was being estimated. Baird's counterexample used all three components and diverged. Baird also proposed a constructive solution, namely residual gradient algorithms that perform gradient descent on the mean-squared Bellman residual, guaranteeing convergence at the cost of a different fixed point.

\subsubsection{Actor-Critic Methods (2000)}

The idea of maintaining both a policy (\textit{actor}) and a value function (\textit{critic}) dates to \citet{barto1983neuronlike}, who used a two-component system to solve the pole-balancing task. Actor-critic methods address the high variance of REINFORCE by replacing Monte Carlo returns with bootstrapped TD targets as the learning signal. The critic learns a value function $V(s)$ by TD updates.
\begin{equation}
V(s_t) \leftarrow V(s_t) + \alpha_c \, \delta_t, \quad \delta_t = r_{t+1} + \gamma V(s_{t+1}) - V(s_t)
\end{equation}
The actor updates the policy using the TD error as a sample of the advantage.
\begin{equation}
\theta \leftarrow \theta + \alpha_a \nabla_\theta \log \pi_\theta(a_t|s_t) \, \delta_t
\end{equation}
The TD error $\delta_t$ estimates $A^\pi(s_t, a_t) = Q^\pi(s_t, a_t) - V^\pi(s_t)$, the advantage of action $a_t$ over the average action.\footnote{That $\delta_t$ is an unbiased estimate of the advantage follows from the policy gradient theorem \citep{SuttonMcAllester2000}.} Positive advantages indicate the action was better than expected; negative advantages indicate it was worse.

\citet{konda2000} provided the first convergence proof for actor-critic algorithms with function approximation, showing convergence to a stationary point under two-timescale learning rates ($\alpha_c \gg \alpha_a$) and a compatibility condition on the critic architecture (Section~\ref{sec:actor_critic}).

%\citet{mnih2016a3c} introduced A3C (Asynchronous Advantage Actor-Critic), which parallelizes actor-critic across multiple CPU threads updating a shared parameter vector. Each thread runs its own copy of the environment, collects trajectories of length $t_{\max}$ or until a terminal state, computes gradients with respect to both actor and critic parameters, and applies updates asynchronously to the shared parameters without synchronization. The gradient for a trajectory $(s_0, a_0, r_0, \ldots, s_{t_{\max}})$ is accumulated by computing $n$-step returns $R_t = \sum_{k=0}^{n-1} \gamma^k r_{t+k} + \gamma^n V(s_{t+n})$ and using these as targets for both policy and value updates. The asynchrony provides implicit exploration, as different threads encounter different states, decorrelating the gradient updates. A3C achieved state-of-the-art results on Atari games while training in half the training time required by DQN on a GPU, using a single multi-core CPU, without the experience replay buffer required by DQN. The synchronous variant, A2C, collects batches from all workers before computing a single update; this sacrifices wall-clock speed for more stable gradients.

\subsubsection{Natural Policy Gradient (2001)}

Standard gradient descent treats all parameter directions equally, but policy parameters define probability distributions whose natural geometry is not Euclidean. \citet{Kakade2001} introduced the natural policy gradient, which measures progress in distribution space rather than parameter space. The update uses the Fisher information matrix $F(\theta)$:
\begin{equation}
F(\theta) = \mathbb{E}_{s \sim d^\pi, a \sim \pi_\theta} \left[ \nabla_\theta \log \pi_\theta(a|s) \nabla_\theta \log \pi_\theta(a|s)^\top \right]
\end{equation}
The natural gradient is
\begin{equation}
\tilde{\nabla}_\theta J(\theta) = F(\theta)^{-1} \nabla_\theta J(\theta)
\end{equation}
This direction is invariant to reparameterization of the policy. In the tabular softmax\footnote{The softmax function maps a vector $\mathbf{z}$ to a probability distribution: $\text{softmax}(z_i) = \exp(z_i)/\sum_j \exp(z_j)$. A softmax policy parameterizes action probabilities as $\pi_\theta(a|s) = \text{softmax}(\theta_{s,a})$.} case, a single natural gradient step with unit step size recovers one step of exact policy iteration (Section~\ref{sec:policy_gradient}).

The computational bottleneck is inverting $F(\theta) \in \mathbb{R}^{d \times d}$. Practical implementations use conjugate gradient methods to solve $F(\theta) x = \nabla_\theta J(\theta)$ without forming $F$ explicitly. This approach was later scaled to deep neural networks by TRPO and PPO.

\subsubsection{Fitted Value Iteration and Fitted Q-Iteration (2005)}
\label{sec:fvi_fqi_algorithms}

Tabular Q-learning maintains a separate entry for every state-action pair. When $|\mathcal{S}|$ is large or the state space is continuous, as in most applied problems, this is infeasible. \emph{Fitted Q-Iteration} (FQI) \citep{Ernst2005} replaces the tabular update with a supervised regression step: given a batch of transitions, fit a function approximator to the Bellman targets.

Let $\mathcal{F}$ be a function class mapping $\mathcal{S} \times \mathcal{A} \to \mathbb{R}$.\label{def:fqi} Initialize $Q_0 \equiv 0$. At each iteration $k = 0, 1, \ldots, K-1$: (i) draw $N$ transitions $(s_i, a_i, r_i, s_i')$ from a generative model; (ii) construct regression targets $y_i^{(k)} = r_i + \gamma \max_{a'} Q_k(s_i', a')$; (iii) set $Q_{k+1} \leftarrow \arg\min_{f \in \mathcal{F}} \frac{1}{N} \sum_{i=1}^N \bigl( f(s_i, a_i) - y_i^{(k)} \bigr)^2$. The output is the greedy policy $\pi_K(s) = \arg\max_a Q_K(s,a)$.

The regression step replaces the exact Bellman application with a projection onto $\mathcal{F}$: $Q_{k+1} = \Pi_{\mathcal{F}} \mathcal{T} Q_k$, where $\mathcal{T}$ is the Bellman optimality operator and $\Pi_{\mathcal{F}}$ is the $L^2$-projection under the sample distribution. Fitted Value Iteration (FVI) \citep{MunosSzepesvari2008} applies the same idea to the value function directly: $V_{k+1} = \Pi_{\mathcal{F}} \mathcal{T}^* V_k$, where $(\mathcal{T}^* V)(s) = \max_a \{ r(s,a) + \gamma \sum_{s'} P(s'|s,a) V(s') \}$.

With feature matrix $\Phi \in \mathbb{R}^{|\mathcal{S}| \times d}$ (rows $\phi(s)^\top$) and per-action weight vectors $\theta_a \in \mathbb{R}^d$, each FQI regression step for action $a$ reduces to the normal equations:
\begin{equation}
\theta_a^{(k+1)} = \bigl(\Phi^\top \Phi\bigr)^{-1} \Phi^\top y_a^{(k)},
\label{eq:fqi_normal}
\end{equation}
where $y_a^{(k)}(s) = r(s,a) + \gamma \sum_{s'} P(s'|s,a) \max_{a'} \phi(s')^\top \theta_{a'}^{(k)}$. Computation is $O(d^2 |\mathcal{S}| + d^3)$ per action per iteration. The FVI update takes the same form with a single weight vector $\theta_V$:
\begin{equation}
\theta_V^{(k+1)} = \bigl(\Phi^\top \Phi\bigr)^{-1} \Phi^\top V_{\mathrm{target}}^{(k)},
\label{eq:fvi_normal}
\end{equation}
where $V_{\mathrm{target}}^{(k)}(s) = \max_a \bigl\{ r(s,a) + \gamma \sum_{s'} P(s'|s,a) \phi(s')^\top \theta_V^{(k)} \bigr\}$. The finite-sample error theory for these methods is developed in Section~\ref{sec:fvi_fqi_theory}.

%%%%%%%%%%%%%%%%%%%%%%%%%%%%%%%%%%%%%%%%%%%%%%%%%%%%%%%%%%%%%%%%%%%%%%%%%%%%%%%
\subsection{The Deep Learning Era}
%%%%%%%%%%%%%%%%%%%%%%%%%%%%%%%%%%%%%%%%%%%%%%%%%%%%%%%%%%%%%%%%%%%%%%%%%%%%%%%

\subsubsection{Deep Q-Networks (2015)}

\citet{mnih2015} trained a single convolutional neural network\footnote{A convolutional neural network applies learned spatial filters to detect local patterns in grid-structured data, commonly used for image inputs.} to play 49 Atari 2600 games directly from pixel inputs ($210 \times 160 \times 3$) and a scalar score, using no game-specific features.

The architecture processed four consecutive frames through three convolutional layers and a fully connected layer.\footnote{A fully connected layer computes $y = Wx + b$ where every input unit is connected to every output unit with learned weights $W$ and biases $b$; it is the affine transformation familiar from linear regression, followed by a nonlinear activation.} The network $Q(s, a; \theta)$ was trained to minimize the squared temporal difference error
\begin{equation}
L(\theta) = \mathbb{E}_{(s,a,r,s') \sim \mathcal{D}} \left[ \left( r + \gamma \max_{a'} Q(s', a'; \theta^-) - Q(s, a; \theta) \right)^2 \right]
\end{equation}

Two innovations stabilized learning. \textit{Experience replay} \citep{Lin1992} stored transitions $(s, a, r, s')$ in a buffer $\mathcal{D}$ and sampled uniformly for training, breaking the temporal correlation between consecutive updates. A \textit{target network} used a frozen copy of parameters $\theta^-$, updated periodically,\footnote{The buffer held $10^6$ transitions; the target network $\theta^-$ was synchronized to $\theta$ every $C = 10{,}000$ steps.} so the regression target does not shift with each gradient step.

%These techniques address the deadly triad identified by Baird. Experience replay creates a more stationary data distribution, partially mitigating off-policy issues. The target network holds the bootstrap target fixed between updates, making the regression problem approximately stationary. Neither technique provides theoretical guarantees, but the empirical success demonstrated that the deadly triad can be managed in practice.

DQN exceeded human-level performance on 29 of 49 games using a single architecture and hyperparameters.\footnote{\emph{Hyperparameters} are design choices fixed before training begins, such as network depth, learning rate, and replay buffer size; they are not estimated by gradient descent.} Games requiring long-horizon planning or sparse rewards, such as Montezuma's Revenge, remained difficult.

\subsubsection{TRPO and PPO (2015, 2017)}

Policy gradient methods suffer from a practical instability: a single large gradient step can move the policy into a region where performance collapses and recovery is slow.

\citet{Schulman2015} addressed this with Trust Region Policy Optimization (TRPO). TRPO solves the constrained optimization problem
\begin{equation}
\max_\theta \; L_{\theta_{\text{old}}}(\theta) \quad \text{subject to} \quad \bar{D}_{\mathrm{KL}}(\theta_{\text{old}}, \theta) \leq \delta
\end{equation}
where $L$ is a surrogate objective based on the advantage function $A^\pi(s,a) = Q^\pi(s,a) - V^\pi(s)$.\footnote{The Kullback-Leibler divergence $D_{\mathrm{KL}}(p \| q) = \sum_x p(x) \log(p(x)/q(x))$ measures statistical distance between distributions. The bar denotes expectation over states: $\bar{D}_{\mathrm{KL}} = \mathbb{E}_{s}[D_{\mathrm{KL}}(\pi_{\theta_{\text{old}}}(\cdot|s) \| \pi_\theta(\cdot|s))]$.}

\citet{Schulman2017} proposed Proximal Policy Optimization (PPO) as a simpler alternative. PPO replaces the KL constraint with a clipped surrogate objective:
\begin{equation}
L^{\text{CLIP}}(\theta) = \mathbb{E}_t \left[ \min\!\left( r_t(\theta) \hat{A}_t, \; \text{clip}(r_t(\theta), 1-\varepsilon, 1+\varepsilon) \hat{A}_t \right) \right]
\end{equation}
where $r_t(\theta) = \pi_\theta(a_t|s_t) / \pi_{\theta_{\text{old}}}(a_t|s_t)$ is the probability ratio. The clipping removes the incentive for $r_t(\theta)$ to move outside the interval $[1-\varepsilon, 1+\varepsilon]$, penalizing large policy updates without explicitly computing a divergence measure.

PPO outperformed A2C, TRPO, and the cross-entropy method on continuous control benchmarks and achieved the highest average reward on 30 of 49 Atari games among the methods tested. PPO is widely used in large-scale RL applications; constraining the magnitude of policy updates is essential for stable optimization.

\subsubsection{Soft Actor-Critic (2018)}

\citet{Haarnoja2018} introduced Soft Actor-Critic (SAC), which adds entropy regularization to the actor-critic framework. The agent maximizes expected return plus an entropy bonus:
\begin{equation}
J(\theta) = \mathbb{E}_{\pi_\theta}\left[\sum_{t=0}^\infty \gamma^t \left(r_t + \tau \mathcal{H}(\pi_\theta(\cdot|s_t))\right)\right]
\end{equation}
where $\mathcal{H}(\pi) = -\sum_a \pi(a) \log \pi(a)$ is the entropy and $\tau > 0$ is a temperature parameter. The entropy bonus encourages exploration by penalizing deterministic policies. The optimal policy under this objective is softmax in the Q-values: $\pi^*(a|s) \propto \exp(Q^*(s,a)/\tau)$, connecting to discrete choice models in econometrics.

SAC maintains two Q-networks (to reduce overestimation bias) and a policy network. The soft Bellman operator for the critic is:
\begin{equation}
(T^\pi Q)(s,a) = r(s,a) + \gamma \mathbb{E}_{s'}\left[ V(s') \right], \quad V(s) = \mathbb{E}_{a \sim \pi}[Q(s,a) - \tau \log \pi(a|s)]
\end{equation}
SAC is off-policy (using experience replay), handles continuous actions naturally, and achieves state-of-the-art sample efficiency on continuous control benchmarks. The entropy regularization provides automatic exploration without $\varepsilon$-greedy schedules.

\FloatBarrier
\subsubsection{Control as Probabilistic Inference}
\label{sec:control_inference}

The \emph{control-as-inference} framework \citep{Todorov2006, Levine2018} recasts RL as probabilistic inference in a graphical model. Appending a binary optimality variable $\mathcal{O}_t$ to the standard MDP, with $P(\mathcal{O}_t = 1 \mid s_t, a_t) = \exp(r(s_t,a_t)/\tau)$, turns policy optimization into posterior inference over actions conditional on future optimality. Structured variational inference (restricting the variational family to the true dynamics) yields the soft Bellman equations already given in the SAC subsection, with the soft value function $V(s) = \tau \log \sum_a \exp(Q(s,a)/\tau)$ identical to McFadden's social surplus function from discrete choice theory \citep{McFadden1978}. Different approximation strategies recover SAC, soft Q-learning, and max-ent policy gradients as special cases.\footnote{The max-ent objective $\max_\pi \sum_t \mathbb{E}[r(s_t,a_t) + \tau \mathcal{H}(\pi(\cdot|s_t))]$ equals the evidence lower bound (ELBO) of the graphical model. Trust-region methods (TRPO, PPO) do not optimize the max-ent objective globally, but each policy update step solves a max-ent subproblem \citep{Schulman2015, Schulman2017}. The $\tau \to 0$ limit recovers standard hard-max algorithms.} Reading the model in reverse, treating the reward as unknown and observed behavior as evidence of optimality, yields maximum entropy inverse reinforcement learning \citep{Ziebart2008}, connecting to the companion survey \citep{RustRawat2026}.

\subsubsection{AlphaGo Zero (2017)}
\label{subsubsec:alphago_zero}

The game of Go has approximately $10^{170}$ legal positions and a branching factor of roughly 250, far beyond brute-force search. \citet{Silver2016} combined reinforcement learning via self-play with Monte Carlo tree search (MCTS) and learned neural networks; the resulting system defeated Lee Sedol in 2016. \citet{Silver2017} then showed that human training data was unnecessary.

AlphaGo Zero uses a single deep residual network $f_\theta(s) = (\mathbf{p}, v)$ that takes a board position $s$ (encoded as raw binary planes without hand-crafted features) and outputs a policy vector $\mathbf{p}$ over legal moves and a scalar value $v$ estimating the probability of winning.\footnote{The input consists of 17 binary planes on the $19 \times 19$ board (eight planes each for Black and White stone positions over the last eight moves, plus one for the current player). The architecture uses 40 residual blocks (79 parameterized layers).}

During play, each move is selected by running MCTS from the current position. Each of 1,600 simulations traverses the search tree by choosing at each node the action maximizing
$$Q(s,a) + c_{\text{puct}} \cdot P(s,a) \cdot \frac{\sqrt{\sum_b N(s,b)}}{1 + N(s,a)}$$
where $Q(s,a)$ is the mean value, $P(s,a)$ is the network's prior, and $N(s,a)$ is the visit count.\footnote{This is a continuous analogue of the upper confidence bound (UCB) strategy, balancing exploitation ($Q$) and exploration ($P/N$).} When the traversal reaches an unexplored position, the neural network evaluates it in a single forward pass; the resulting value $v$ propagates back up the path, updating $Q$ and $N$ at each edge. The move with the highest root visit count is selected.


The training loop generates self-play games. At each position $s_t$, MCTS produces improved move probabilities $\boldsymbol{\pi}_t$ (proportional to visit counts). At game end, the outcome $z \in \{-1, +1\}$ is recorded. Each position yields a training triple $(s_t, \boldsymbol{\pi}_t, z)$, and the network minimizes
\begin{equation}
\ell(\theta) = (z - v)^2 - \boldsymbol{\pi}^\top \log \mathbf{p} + c \|\theta\|^2
\end{equation}
combining value prediction loss, policy cross-entropy loss, and $L_2$ regularization. MCTS acts as a policy improvement operator, since the search probabilities $\boldsymbol{\pi}$ are stronger than the raw network outputs $\mathbf{p}$; training the network to match $\boldsymbol{\pi}$ distills search improvements back into the network.\footnote{The system that defeated Lee Sedol used fixed network weights throughout the match; the months of self-play constituted the training phase, while the five-game match was purely execution (Section~\ref{section:language}).}

Go was well-suited to this architecture. Its fixed $19 \times 19$ board maps naturally to convolutional networks, its perfect information and deterministic transitions make MCTS's tree structure exact, and the binary game outcome provides an unambiguous training signal. \citet{Igami2020} interprets the architecture in econometric terms, where the policy network is a conditional choice probability (CCP) estimator, the value network is a conditional value function (CVF) estimator, and the system performs CCP estimation and forward simulation jointly, connecting to the approach of \citet{HotzMiller1993} in dynamic discrete choice.
